\begin{theorem}[zeta mul zeta]
\label{zeta_mul_zeta}
\lean{ArithmeticFunction.zeta_mul_zeta}
% at /home/runner/work/trough/trough/library/PrimeNumberTheoremAnd/IwaniecKowalskiCh1.lean:147.0-159.29
\leanok
The Dirichlet convolution of $\zeta$ with itself is $\tau$ (the divisor count function).
\end{theorem}
\begin{proof}
\leanok
By definition of $\zeta$, we have $\zeta(n) = 1$ for all $n \geq 1$. Thus, the Dirichlet convolution
  $(\zeta * \zeta)(n)$ counts the number of ways to write $n$ as a product of two positive integers,
  which is exactly the number of divisors of $n$, i.e., $\tau(n)$.
\end{proof}
